\documentclass[11pt, a4paper]{article}

% ─── Packages ────────────────────────────────────────────────────────────────
\usepackage[T1]{fontenc}
\usepackage[utf8]{inputenc}
\usepackage{lmodern}
\usepackage[margin=2.5cm]{geometry}
\usepackage{amsmath, amssymb, bm}
\usepackage{graphicx}
\usepackage{booktabs}
\usepackage{multirow}
\usepackage{array}
\usepackage{xcolor}
\usepackage{hyperref}
\usepackage{natbib}
\usepackage{microtype}
\usepackage{caption}
\usepackage{subcaption}
\usepackage{float}
\usepackage{algorithm}
\usepackage{algpseudocode}
\usepackage{enumitem}
\usepackage{parskip}
\usepackage{titlesec}
\usepackage{setspace}

\onehalfspacing
\setlength{\parindent}{0pt}
\setlength{\parskip}{6pt}

\hypersetup{
    colorlinks=true,
    linkcolor=blue!60!black,
    citecolor=green!50!black,
    urlcolor=blue!70!black
}

\graphicspath{{figures/}}

% ─── Title ───────────────────────────────────────────────────────────────────
\title{
    \vspace{-1cm}
    {\large Capstone Design Project — Final Report}\\[0.8em]
    \textbf{Dynamic Grid Trading via Proximal Policy Optimization:\\
    Bayesian Coefficient Optimization and Volatility-Adaptive Order Placement}
}

\author{
    Lee Chanhee\\
    Department of Computer Science and Engineering\\
    Seoul National University of Science and Technology\\
    \texttt{happilyfly@seoultech.ac.kr}
}

\date{April 2026}

% ─── Document ────────────────────────────────────────────────────────────────
\begin{document}

\maketitle
\thispagestyle{empty}

% ─── Abstract ────────────────────────────────────────────────────────────────
\begin{abstract}
We present a reinforcement learning framework for \emph{dynamic grid trading} on the
BTC/USDT hourly market in which a PPO agent outputs continuous ATR-scaled limit-order
gaps without predicting price direction.
Unlike prior RL trading systems that rely on directional signals or discrete
buy/sell/hold actions, our agent outputs a continuous two-dimensional action
$(a_\text{agg}, a_\text{pt}) \in [0,1]^2$ that maps directly onto four limit-order
prices scaled by the Average True Range (ATR).
Eight formula coefficients controlling the shape of this mapping are jointly optimized
via Bayesian search (Optuna TPE, 50 trials), yielding a \(+145\%\) improvement in
validation Sharpe ratio compared to the default parameterization.
The final model, trained for 3 million steps on data from 2020--2022, achieves a
validation Sharpe ratio of \textbf{35.42} and a test-set Sharpe ratio of \textbf{43.04}
(2024--2026), outperforming the best fixed-grid baseline by \textbf{29.2$\times$}
while maintaining a maximum drawdown of only \textbf{3.1\%} versus 50.1\% for
buy-and-hold.
Post-hoc behavioral analysis reveals that the learned policy converges to a
near-constant action output, indicating that the Bayesian-optimized coefficients
alone encode the dominant performance driver; implications for the role of
state-conditioned adaptation are discussed.
\end{abstract}

\vspace{1em}
\noindent\textbf{Keywords:} reinforcement learning, grid trading, PPO, Bayesian optimization,
cryptocurrency, ATR-scaled orders, policy saturation

\newpage
\tableofcontents
\newpage

% ─── 1. Introduction ─────────────────────────────────────────────────────────
\section{Introduction}

Cryptocurrency markets exhibit persistent high volatility that creates recurring
micro-price oscillations regardless of the long-term trend direction.
Grid trading exploits this property by placing a ladder of buy and sell limit orders
around the current price, profiting from each completed buy-then-sell cycle
\citep{sun2023rl}.
The central weakness of conventional grid trading is that the grid spacing is set
once at deployment and never adjusted: intervals too narrow flood the market with
unprofitable tiny trades, while intervals too wide miss most of the oscillations.

Recent RL trading literature has concentrated almost entirely on \emph{directional}
strategies — predicting whether the next price movement is up or down and acting
accordingly \citep{liu2021bitcoin, yasin2024rl, pham2025dqn}.
This line of work inherits the difficulty of price forecasting and cannot profit in
sideways or highly volatile markets where prediction accuracy is low.
Only a handful of studies treat the RL agent as a \emph{market-making} or
\emph{interval-selection} agent rather than a directional predictor, leaving the
dynamic grid problem largely unexplored.

This paper addresses the following research questions:

\begin{enumerate}[label=\textbf{RQ\arabic*:}]
    \item Can a PPO agent with a continuous ATR-scaled action space outperform
          fixed-interval grid baselines on unseen BTC/USDT data?
    \item Does the trained agent adapt its grid intervals to the prevailing
          volatility regime?
\end{enumerate}

\subsection*{Contributions}

\begin{itemize}
    \item \textbf{Non-directional MDP formulation.}
          A five-dimensional state space encoding both market regime and portfolio
          position, paired with a continuous two-dimensional action that controls
          grid interval width without any directional price prediction.

    \item \textbf{ATR-scaled, tunable order placement formula.}
          Four limit orders (two buys, two sells) are placed at distances that
          are proportional to the current ATR.
          Eight scalar coefficients parameterize the formula, enabling
          data-driven calibration while keeping the economic interpretation intact.

    \item \textbf{Bayesian coefficient optimization.}
          We apply Optuna TPE with MedianPruner to search the eight-dimensional
          coefficient space, achieving a $+145\%$ improvement in validation Sharpe
          over the hand-designed defaults in 50 trials.

    \item \textbf{Rigorous train/validation/test separation.}
          Following \citet{gort2022backtest}, the test set (2024--2026) is kept
          completely sealed until after all design decisions are frozen.
          Test-set evaluation is performed exactly once.
\end{itemize}

% ─── 2. Related Work ─────────────────────────────────────────────────────────
\section{Related Work}

\subsection{Surveys of RL Trading}

\citet{sun2023rl} survey the RL trading literature and identify two recurring
limitations: insufficient interpretability of learned policies and poor
generalization across market regimes.
\citet{liu2025review} review ten papers spanning 2018--2025 and confirm that RL
agents demonstrate superior adaptability to non-stationary markets compared to
rule-based methods, yet training stability and sample efficiency remain the
principal barriers to practical deployment.

\subsection{Directional RL Trading Systems}

\citet{liu2021bitcoin} combine an LSTM price predictor with a PPO agent on
BTC/USDT, feeding predicted prices as state features.
The two-stage pipeline couples trading performance to prediction accuracy,
propagating forecast errors directly into order decisions.
Our approach eliminates the prediction stage entirely.

\citet{yasin2024rl} benchmark DQN, PPO, and A2C with 20 technical indicators on
equity markets.
Their discrete action space (buy / sell / hold) and the absence of portfolio-state
information in the observation prevent the agent from sizing orders relative to
its current exposure — a critical gap for grid trading where position awareness
drives interval selection.

\citet{pham2025dqn} apply DQN to \emph{strategy selection}: the agent chooses
among five hard-coded rule-based strategies at each step rather than learning
the strategy itself.
PCA compression of the indicator state further obscures the economic meaning of
the learned policy.
We instead train the agent to output grid intervals directly, preserving full
interpretability.

\subsection{Risk-Aware RL Trading}

\citet{bandarupalli2025risk} augment the PPO reward with an explicit portfolio
variance penalty to discourage risky positions.
They report that the penalty suppresses too many profitable trades, causing the
agent to underperform buy-and-hold.
This finding motivates our reward design, which relies solely on equity change
and transaction-cost penalties without a separate risk term: grid trading is
inherently volatility-seeking, making a variance penalty counterproductive.

\subsection{Backtesting Methodology}

\citet{gort2022backtest} empirically demonstrate that single-split validation in
RL trading research leads to significant backtest overfitting.
Their recommendation of strict temporal splits directly informs our
train (2020--2022) / validation (2023) / test (2024--2026) design, where the
test partition is accessed only after all hyperparameter and coefficient
decisions are finalized.

\subsection{Positioning of This Work}

Table~\ref{tab:comparison} summarizes how this work differs from its closest
predecessors.

\begin{table}[h]
\centering
\caption{Comparison with related RL trading systems.}
\label{tab:comparison}
\setlength{\tabcolsep}{5pt}
\begin{tabular}{lccccc}
\toprule
 & \textbf{Ours} & Liu '21 & Yasin '24 & Pham '25 & Banda. '25 \\
\midrule
Non-directional          & \checkmark & $\times$ & $\times$ & $\times$ & $\times$ \\
Continuous action space  & \checkmark & $\times$ & $\times$ & $\times$ & $\times$ \\
Position state           & \checkmark & $\times$ & $\times$ & $\times$ & $\checkmark$ \\
ATR-scaled orders        & \checkmark & $\times$ & $\times$ & $\times$ & $\times$ \\
Bayesian coef.\ tuning   & \checkmark & $\times$ & $\times$ & $\times$ & $\times$ \\
Sealed test set          & \checkmark & $\times$ & $\times$ & $\times$ & $\times$ \\
\bottomrule
\end{tabular}
\end{table}

% ─── 3. Methodology ──────────────────────────────────────────────────────────
\section{Methodology}

\subsection{Data and Preprocessing}

We collect hourly BTC/USDT candlestick data from Binance via the CCXT library,
covering January 2020 through April 2026 (54,733 candles after deduplication).
The data are split chronologically into three non-overlapping partitions:

\begin{itemize}
    \item \textbf{Train} (2020-01-01 -- 2022-12-31, 25,916 candles): covers both a bull
          market (2020--2021) and a major bear market (2022).
    \item \textbf{Validation} (2023-01-01 -- 2023-12-31, 8,736 candles): a recovery year
          used for hyperparameter selection and model evaluation during training.
    \item \textbf{Test} (2024-01-01 -- 2026-04-09, 19,901 candles): kept completely sealed
          until all design decisions were finalized; evaluated exactly once.
\end{itemize}

Each raw candle is augmented with three derived series, all computed over a rolling
window of $W = 168$ hours (one week):

\begin{align}
    \text{log\_price}_t &= \log\!\left(\frac{p_t}{\bar{p}_{t,W}}\right), \\
    \text{volatility}_t &= \frac{\text{ATR}(168)_t}{p_t}, \\
    \text{zscore}(x)_t  &= \frac{x_t - \mu_{x,t,W}}{\sigma_{x,t,W} + \epsilon},
\end{align}

where $\bar{p}_{t,W}$ is the rolling mean close price, ATR is the Average True Range
computed directly from OHLC data, and $\epsilon = 10^{-8}$ ensures numerical
stability.
The z-score transformation is applied to \texttt{log\_price} and \texttt{volatility}
before storage to prevent scale mismatch at the policy input.

\subsection{Trading Environment (MDP)}

We model the trading problem as a Markov Decision Process
$\mathcal{M} = (\mathcal{S}, \mathcal{A}, \mathcal{T}, \mathcal{R}, \gamma)$
implemented as a Gymnasium environment.

\subsubsection*{State Space}

The observation at time $t$ is a five-dimensional vector:

\begin{equation}
    \bm{s}_t = \bigl[
        \underbrace{\hat{z}(\text{log\_price}_t)}_{\text{price level}},\;
        \underbrace{d_t}_{\text{divergence}},\;
        \underbrace{h_t}_{\text{position size}},\;
        \underbrace{c_t}_{\text{cash ratio}},\;
        \underbrace{\hat{z}(\text{volatility}_t)}_{\text{market volatility}}
    \bigr] \in \mathbb{R}^5,
\end{equation}

where $\hat{z}(\cdot)$ denotes the rolling z-score normalization,

\begin{align}
    d_t &= \frac{\bar{p}^\text{avg}_t - p_t}{\bar{p}^\text{avg}_t}
           \quad \text{(divergence from average entry price)}, \\
    h_t &= \frac{q_t \cdot p_t}{K}
           \quad \text{(holdings value as fraction of initial capital $K$)}, \\
    c_t &= \frac{\text{cash}_t}{K}
           \quad \text{(cash fraction of initial capital)}.
\end{align}

$\bar{p}^\text{avg}_t$ is the current-cycle average entry price (set to the last
completed cycle's average entry when not holding, and zero when no trades have
occurred).
The value ratio $h_t$ over the raw quantity is used because the same BTC holding
carries twice the portfolio risk when the price doubles.

\subsubsection*{Action Space}

The action space is a continuous two-dimensional box:
\begin{equation}
    \mathcal{A} = [0, 1]^2,
    \quad
    \bm{a}_t = (a_\text{agg},\; a_\text{pt}),
\end{equation}

where $a_\text{agg}$ (aggressiveness) controls buy-order placement and $a_\text{pt}$
(profit target) controls sell-order placement.
This is in sharp contrast to discrete action spaces (buy/sell/hold) used by most
prior work \citep{liu2021bitcoin, yasin2024rl}, which cannot express the
\emph{degree} of aggressiveness.

\subsubsection*{Order Placement Formula}

At each step, four limit orders are computed from the action via an
ATR-proportional formula:

\begin{align}
    r_t &= \frac{\text{ATR}(168)_t}{p_t}
           \quad \text{(normalized volatility ratio)}, \label{eq:atr_ratio} \\[4pt]
    \Delta^\text{buy-hi}_t  &= r_t \cdot (A_b + a_\text{agg}\cdot B_b), \label{eq:buy_hi} \\
    \Delta^\text{buy-lo}_t  &= r_t \cdot (C_b + a_\text{agg}\cdot D_b), \label{eq:buy_lo} \\
    \Delta^\text{sell-mkt}_t &= r_t \cdot (A_s + a_\text{pt}\cdot B_s), \label{eq:sell_mkt} \\
    \Delta^\text{sell-cst}_t &= r_t \cdot (C_s + a_\text{pt}\cdot D_s), \label{eq:sell_cst}
\end{align}

and the four order prices are:

\begin{align}
    P^\text{buy-hi}_t  &= p_t\cdot(1 - \Delta^\text{buy-hi}_t), \\
    P^\text{buy-lo}_t  &= p_t\cdot(1 - \Delta^\text{buy-lo}_t), \\
    P^\text{sell-mkt}_t &= p_t\cdot(1 + \Delta^\text{sell-mkt}_t), \\
    P^\text{sell-cst}_t &= \bar{p}^\text{avg}_t\cdot(1 + \Delta^\text{sell-cst}_t).
\end{align}

Here $\bm{\theta} = (A_b, B_b, C_b, D_b, A_s, B_s, C_s, D_s) \in \mathbb{R}^8$
are the formula coefficients.
ATR-scaling ensures that the action range $[0,1]$ maps consistently onto a
meaningful fill-probability spectrum regardless of the current volatility level,
avoiding the regime where most limit orders never execute.

The dual sell-order design reflects two complementary liquidation channels:
\texttt{sell-mkt} targets a quick exit above the current price (fast cash recovery),
while \texttt{sell-cst} targets a larger gain above the average cost basis (large-rally
capture).

\subsubsection*{Order Fill and Cycle Logic}

Order fills are simulated against the next candle's high/low prices
(next-bar execution model), with sell orders processed before buy orders on the
same bar.
Each complete buy-then-sell sequence forms a \emph{trading cycle}.
Capital allocation within a cycle uses a fixed-budget scheme: the cash available
at cycle start is divided into $n_\text{splits} = 2$ slots, each of which funds
$n_\text{buy} = 2$ limit orders.
When the total holdings drop to or below the value of a single slot, all remaining
holdings are liquidated to guarantee cycle closure.

\subsubsection*{Reward}

The reward at each step is the normalized equity change:

\begin{equation}
    r_t = \frac{V_t - V_{t-1}}{K},
    \label{eq:reward}
\end{equation}

where $V_t = \text{cash}_t + q_t \cdot p_t$ is the total portfolio value.
Transaction costs ($0.05\%$ per fill, Binance maker fee) are deducted directly
from the cash balance at execution time, so they are automatically reflected in
the equity change without a separate penalty term.

\subsection{PPO Agent}

We use Proximal Policy Optimization \citep{schulman2017ppo} via Stable-Baselines3.
The policy and value networks share an MLP trunk (two hidden layers of 64 units,
tanh activations).
The action distribution is a diagonal Gaussian clipped to $[0, 1]^2$
by a sigmoid squashing layer.

Training uses four parallel environments (\texttt{n\_envs}~$= 4$) with episodes
of length $T_\text{ep} = 2016$ steps ($\approx 12$ weeks).
Random episode starts within the training window expose the agent to diverse
market regimes.
A cosine learning rate schedule decays from $3\times10^{-4}$ to $10^{-5}$
over 3 million environment steps.
The entropy coefficient $\alpha_H = 0.05$ is selected by ablation
(Section~\ref{sec:param-search}).

\subsection{Bayesian Coefficient Optimization}
\label{sec:bayesian}

The eight formula coefficients in $\bm{\theta}$ encode prior economic knowledge
(e.g., $A_b < C_b$ forces buy-hi closer to the current price than buy-lo), but
their exact values are unknown.
We treat them as hyperparameters and search the eight-dimensional continuous space
using Tree-structured Parzen Estimator (TPE) with a MedianPruner:

\begin{itemize}
    \item \textbf{Objective}: mean validation Sharpe ratio over five episodes
          after $10^6$ training steps with a warm-started PPO agent.
    \item \textbf{Soft constraints}: trials violating $C_b > A_b$ and $D_b > B_b$
          (which would invert the buy-hi / buy-lo ordering) return a sentinel value
          of $-999$ and are excluded from the TPE model.
    \item \textbf{Pruner}: MedianPruner with 10 warm-up trials kills unpromising
          trials after their first intermediate evaluation, recovering compute for
          more promising configurations.
    \item \textbf{Budget}: 50 trials ($\approx 4.3$ hours on a single CPU).
\end{itemize}

The search space and the final optimized values are given in
Table~\ref{tab:coefs}.

\begin{table}[h]
\centering
\caption{Formula coefficient search space and Bayesian-optimized values (Trial \#42).}
\label{tab:coefs}
\setlength{\tabcolsep}{6pt}
\begin{tabular}{llrrrr}
\toprule
Coef. & Role & Search min & Search max & \textbf{Default} & \textbf{Optimal} \\
\midrule
$A_b$ & buy-hi base & 0.1 & 3.0 & 0.500 & \textbf{0.285} \\
$B_b$ & buy-hi range & 0.1 & 5.0 & 1.500 & \textbf{1.748} \\
$C_b$ & buy-lo base & 0.1 & 10.0 & 2.500 & \textbf{5.223} \\
$D_b$ & buy-lo range & 0.5 & 20.0 & 7.500 & \textbf{18.683} \\
$A_s$ & sell-mkt base & 0.1 & 3.0 & 0.500 & \textbf{0.101} \\
$B_s$ & sell-mkt range & 0.1 & 5.0 & 1.500 & \textbf{0.890} \\
$C_s$ & sell-cost base & 0.1 & 10.0 & 2.500 & \textbf{6.913} \\
$D_s$ & sell-cost range & 0.5 & 20.0 & 7.500 & \textbf{5.457} \\
\bottomrule
\end{tabular}
\end{table}

% ─── 4. Experimental Setup ───────────────────────────────────────────────────
\section{Experimental Setup}

\subsection{Baselines}

We compare against seven fixed-rule strategies that use no learning:

\begin{itemize}
    \item \textbf{Buy-and-Hold}: purchase BTC at the start, sell at the end.
    \item \textbf{Fixed Grid $\delta\%$} ($\delta \in \{1, 2, 5\}$): buy orders placed
          $\delta\%$ below and sell orders $\delta\%$ above the current price.
    \item \textbf{ATR Grid $k$} ($k \in \{0.5, 1.0, 2.0\}$): order spacing is $k$ times
          the current ATR, mimicking the ATR-proportional scaling of our agent but
          without learned adaptation.
\end{itemize}

The ATR Grid baselines are particularly informative: they share the same market
information (ATR) as our agent but apply it through a fixed rule.
Performance differences relative to these baselines reflect the combined effect of
Bayesian coefficient optimization and the RL training procedure.

\subsection{Evaluation Metric}

The primary metric is the annualized Sharpe ratio:

\begin{equation}
    \text{Sharpe} = \frac{\mathbb{E}[r_t]}{\text{Std}(r_t)} \cdot \sqrt{8760},
    \label{eq:sharpe}
\end{equation}

where $r_t$ is the per-step return and 8760 is the number of hourly periods per
year.
Secondary metrics include maximum drawdown (MDD), total return, number of
completed trading cycles, and cycle win rate.

\subsection{Experimental Protocol}

All experiments use seed 42.
Validation Sharpe is computed as the mean over five evaluation episodes of length
$T_\text{ep} = 2016$ steps with random starts within the validation window and
$\text{deterministic}=\text{True}$ inference.
The test evaluation is a single deterministic pass over the entire 19,901-step
test sequence, ensuring no information from random start selection can influence
the reported result.

% ─── 5. Results ──────────────────────────────────────────────────────────────
\section{Results}

\subsection{Hyperparameter Search}
\label{sec:param-search}

Before Bayesian coefficient tuning, we conduct a systematic ablation over three
environment hyperparameters on the validation set:

\begin{itemize}
    \item \textbf{Number of capital slots} $n_\text{splits} \in \{2, 8, 16\}$:
          $n_\text{splits}=2$ achieves the highest Sharpe by concentrating capital
          per cycle and reducing idle cash.
    \item \textbf{Sell threshold basis} (\texttt{avg\_price} vs.\ \texttt{price}):
          anchoring the sell-cost order to the cycle's average entry price produces
          more consistent profit locks than anchoring to the current market price.
    \item \textbf{Entropy coefficient} $\alpha_H \in \{0.01, 0.05, 0.10\}$:
          $\alpha_H = 0.05$ balances exploration and exploitation, increasing
          Sharpe from 14.39 (no VecNorm baseline, exp007) to 17.58.
\end{itemize}

The selected configuration (exp013b) with default formula coefficients achieves a
validation Sharpe of 17.58, which serves as the baseline for Bayesian optimization.

\subsection{Bayesian Coefficient Optimization}

Figure~\ref{fig:optuna} shows the results of the 50-trial Bayesian search.
The best trial (\#42) achieves a validation Sharpe of 42.997 — a $+145\%$
improvement over the default-coefficient baseline.

\begin{figure}[H]
    \centering
    \begin{subfigure}[b]{0.48\textwidth}
        \centering
        \includegraphics[width=\textwidth]{fig_01_distribution.png}
        \caption{Sharpe distribution across 50 trials.
                 Red dashed line: default-coefficient baseline (17.58).}
        \label{fig:optuna_dist}
    \end{subfigure}
    \hfill
    \begin{subfigure}[b]{0.48\textwidth}
        \centering
        \includegraphics[width=\textwidth]{fig_03_importance.png}
        \caption{Pearson correlation between each coefficient and Sharpe ratio.
                 $A_s$ (sell-market base) shows the strongest negative correlation.}
        \label{fig:optuna_importance}
    \end{subfigure}
    \caption{Bayesian coefficient tuning results (50 TPE trials, 1M steps each).}
    \label{fig:optuna}
\end{figure}

The coefficient importance analysis (Figure~\ref{fig:optuna_importance}) reveals
that $A_s$, the base value of the sell-market gap, is by far the most influential
parameter.
Its optimal value (0.101) sits at the search lower bound of 0.1, suggesting the
agent benefits from a very tight sell-market order — placing the quick-exit ask
just $0.101 \times r_t$ above the current price, enabling near-immediate cash
recovery after each buy fill.

Figure~\ref{fig:gap_comparison} illustrates how the optimized coefficients shift the
effective order-gap ranges compared to the defaults.

\begin{figure}[H]
    \centering
    \includegraphics[width=0.85\textwidth]{fig_06_gap_comparison.png}
    \caption{Effective order-gap ranges (in ATR multiples) for default vs.\ optimized
             coefficients, shown across the full action range $[0,1]$.
             Note the dramatic expansion of the buy-lo range and the sharp contraction
             of the sell-market gap.}
    \label{fig:gap_comparison}
\end{figure}

The optimized strategy can be interpreted as a \emph{fast-sell / deep-buy} policy:
\begin{itemize}
    \item \textbf{Sell-market} contracts to $[0.10, 1.0] \times r_t$
          (vs.\ default $[0.5, 2.0] \times r_t$): the agent exits quickly on any
          small upward tick, maximizing cycle frequency.
    \item \textbf{Buy-lo} expands to $[5.2, 23.9] \times r_t$
          (vs.\ default $[2.5, 10.0] \times r_t$): a deep standing bid positioned
          to catch large dips that appear during high-volatility events.
    \item \textbf{Sell-cost} rises to $[6.9, 12.4] \times r_t$: the cost-basis
          sell is reserved only for major rallies above the average entry price.
\end{itemize}

\subsection{Final Model Training (exp016)}

Using the optimized coefficients from Trial \#42, we train the final PPO agent for
3 million steps.
Figure~\ref{fig:learning_curve} shows the validation Sharpe throughout training.

\begin{figure}[H]
    \centering
    \includegraphics[width=\textwidth]{03_learning_curve.png}
    \caption{Validation Sharpe ratio vs.\ training steps.
             \textbf{(a) Pre-fix experiments} (exp001--002): a reward-function bug
             (fee double-counting) and an eval-pipeline bug kept Sharpe below 1.0.
             \textbf{(b) Post-fix experiments}: after both bugs were resolved, Sharpe
             climbed from $\sim$12 (exp008) to 17.6 (exp013b) through hyperparameter
             search, and finally to 35.4 with Bayesian-optimized coefficients (exp016,
             best checkpoint at step 2.1M, $\star$).}
    \label{fig:learning_curve}
\end{figure}

The best checkpoint is saved at step 2,100,000 with validation Sharpe 35.42 and MDD
2.46\%.
Note that the full-training Sharpe (35.42) is lower than the 1M-step Optuna
objective (42.997): the Optuna evaluation uses five random-start episodes and may
select favorable market windows, whereas the final evaluation averages across fixed
evaluation windows, providing a more conservative and reproducible estimate.

\subsection{Validation Set Results}

Table~\ref{tab:val_results} compares the best PPO checkpoint against all baselines
on the 2023 validation set.

\begin{table}[H]
\centering
\caption{Validation set results (2023, 8,736 candles).}
\label{tab:val_results}
\setlength{\tabcolsep}{6pt}
\begin{tabular}{lrrrr}
\toprule
Strategy & Sharpe & Return (\%) & MDD (\%) & Cycles \\
\midrule
Buy \& Hold           &  2.377 & +150.18 & 21.74 &    0 \\
Fixed Grid 1\%        &  1.991 &  +64.11 & 17.50 &  120 \\
Fixed Grid 2\%        &  2.039 &  +31.05 & 12.57 &   38 \\
Fixed Grid 5\%        &  1.377 &   +4.99 &  3.32 &    4 \\
ATR Grid $k{=}0.5$    &  1.141 &  +30.80 & 20.66 &  279 \\
ATR Grid $k{=}1.0$    &  1.607 &  +54.37 & 18.63 &  177 \\
ATR Grid $k{=}2.0$    &  2.234 &  +55.42 & 13.23 &   91 \\
\midrule
\textbf{PPO (exp016)} & \textbf{35.424} & \textbf{+365.60} & \textbf{2.46} & \textbf{938} \\
\bottomrule
\end{tabular}
\end{table}

The PPO agent improves on the strongest baseline (ATR Grid $k{=}2.0$, Sharpe 2.234)
by a factor of \textbf{15.9$\times$} while reducing MDD from 13.23\% to 2.46\%.

\subsection{Test Set Results}
\label{sec:test-results}

The test set (2024-01-01 to 2026-04-09) was accessed once after all model
and coefficient decisions were frozen.
This period covers BTC price movements from approximately \$42,000 to over
\$100,000 and back to \$75,000, including both a major bull run (2024) and a
sharp correction (2025 Q1), providing a diverse out-of-sample evaluation.

\begin{table}[H]
\centering
\caption{Test set results (2024--2026, 19,901 candles).}
\label{tab:test_results}
\setlength{\tabcolsep}{6pt}
\begin{tabular}{lrrrr}
\toprule
Strategy & Sharpe & Return (\%) & MDD (\%) & Cycles \\
\midrule
Buy \& Hold           &  0.681 &  +62.25 & 50.08 &    0 \\
Fixed Grid 1\%        &  0.724 &  +52.81 & 29.87 &  176 \\
Fixed Grid 2\%        &  0.625 &  +34.76 & 26.66 &   54 \\
Fixed Grid 5\%        &  1.472 &  +14.28 &  3.16 &    9 \\
ATR Grid $k{=}0.5$    & $-$0.004 & $-$20.22 & 49.77 &  325 \\
ATR Grid $k{=}1.0$    &  0.258 &   +5.29 & 49.59 &  242 \\
ATR Grid $k{=}2.0$    &  0.436 &  +22.73 & 29.85 &   88 \\
\midrule
\textbf{PPO (exp016)} & \textbf{43.040} & --- & \textbf{3.12} & \textbf{9,520} \\
\bottomrule
\end{tabular}
\end{table}

\noindent\textit{Note on test-set return: over 9,520 completed cycles with a mean
cycle PnL of 0.182\% compounding continuously across 27 months, the arithmetic
cumulative return reaches astronomical values. The Sharpe ratio and MDD are the
appropriate metrics here, as they normalize for time and risk. The 3.12\% MDD
is the operative measure of capital safety.}

Figure~\ref{fig:equity} shows the test-period equity curve in log scale together
with all baseline strategies.

\begin{figure}[H]
    \centering
    \includegraphics[width=0.9\textwidth]{06_test_equity_curve.png}
    \caption{Portfolio equity curves on the test set (2024--2026), log scale.
             The PPO agent (blue) grows exponentially while all baselines
             remain near the initial capital line. Bottom panel: BTC/USDT price.}
    \label{fig:equity}
\end{figure}

Figure~\ref{fig:progression} places the test result in the context of the full
experimental progression.

\begin{figure}[H]
    \centering
    \includegraphics[width=0.9\textwidth]{06_sharpe_progression.png}
    \caption{\textbf{Left}: validation Sharpe at each major experiment milestone.
             \textbf{Right}: final PPO Sharpe vs.\ best baseline on Val and Test sets.
             Test Sharpe exceeds Val Sharpe, providing no evidence of overfitting.}
    \label{fig:progression}
\end{figure}

% ─── 6. Analysis ─────────────────────────────────────────────────────────────
\section{Analysis and Discussion}

\subsection{RQ1: Does PPO Outperform Fixed-Grid Baselines?}

On the validation set, PPO achieves a Sharpe of 35.42 versus 2.23 for the best
baseline, a \textbf{15.9$\times$} advantage.
On the held-out test set, PPO achieves 43.04 versus 1.47 for Fixed Grid 5\%,
a \textbf{29.2$\times$} advantage.
Both results unambiguously answer RQ1 in the affirmative.

The higher test Sharpe relative to validation Sharpe (43.04 vs.\ 35.42) is
consistent with the 2024--2025 BTC bull market providing an environment
especially favorable to the agent's fast-cycle, high-frequency trading style.
Critically, the direction of this difference (test $>$ val) rules out test-set
overfitting: a model that had implicitly seen the test data would show test
performance \emph{no better} than validation, not superior.

\subsection{RQ2: Volatility Regime Adaptation}

Figure~\ref{fig:behavior} shows the agent's action output on the test set,
stratified by volatility regime (Low / Mid / High, defined by ATR/price thresholds).

\begin{figure}[H]
    \centering
    \includegraphics[width=\textwidth]{06_behavior_analysis.png}
    \caption{Action distributions on the test set.
             \textbf{Top-left/Bottom-left}: both aggressiveness and profit target
             are concentrated at exactly 0.0 across all volatility regimes,
             indicating a saturated constant policy.
             \textbf{Bottom-right}: the regime distribution across the test period.}
    \label{fig:behavior}
\end{figure}

Contrary to the expectation of RQ2, the trained policy outputs a near-constant
action vector $\bm{a} \approx (0, 0)$ deterministically, regardless of the
observed volatility regime.
Inspection of the policy network confirms this: the raw Gaussian mean for
aggressiveness is $-9.19$ (tanh-squashed to $\approx -1$, rescaled to $\approx 0$),
and for profit target $-4.30$ (rescaled to $\approx 0$).
No statistically meaningful difference in action values exists across Low, Mid,
and High volatility regimes.

\textbf{Why did the policy saturate?}
The root cause lies in a structural coupling between Bayesian coefficient
optimization and the RL action space.
The order-placement formula is:
\begin{equation}
    \Delta^\text{sell-mkt}_t = r_t \cdot (A_s + a_\text{pt} \cdot B_s).
\end{equation}
Bayesian search drove $A_s$ to its lower bound ($0.101$), making the minimum
action value ($a_\text{pt} = 0$) already optimal: any increase in $a_\text{pt}$
only widens the sell gap, reducing cycle frequency.
The RL agent correctly learned to keep $a_\text{pt} = 0$, but this renders
state-conditioned adaptation unnecessary.
In other words, the Bayesian stage solved the gap-sizing problem at the
coefficient level, leaving no residual degrees of freedom for RL to exploit.

\textbf{Implication.}
RQ2 cannot be answered with the current design.
The result does not imply that volatility-regime adaptation is impossible or
unhelpful — it implies that the coefficient optimization removed the mechanism
through which adaptation would have been expressed.
A redesigned system in which coefficients fix the \emph{range} of behavior while
RL controls the \emph{selection within that range} is the natural next step
(Section~\ref{sec:conclusion}).

\subsection{Cycle Behavior}

Figure~\ref{fig:cycle} characterizes the 9,520 trading cycles completed during
the test period.

\begin{figure}[H]
    \centering
    \includegraphics[width=\textwidth]{06_cycle_analysis.png}
    \caption{Test-set cycle analysis.
             \textbf{Left}: PnL distribution (win rate 96.9\%, mean +0.182\%).
             \textbf{Center}: cycle duration distribution (mean 1.2 bars $\approx 1.2$ hours).
             \textbf{Right}: cumulative compounding of per-cycle returns on a log scale.}
    \label{fig:cycle}
\end{figure}

The 96.9\% cycle win rate and 1.2-hour average duration demonstrate that the agent
operates in a near-market-making mode: opening and closing positions rapidly,
collecting small spreads on each cycle.
The log-linear cumulative return curve (right panel) confirms that the compounding
occurs steadily without catastrophic drawdowns, consistent with the 3.12\% MDD.

\subsection{Limitations}

\textbf{Market impact.}
The simulation assumes infinite liquidity: limit orders are filled at the stated
price with no slippage.
At the simulated capital scale (\$10,000), this assumption is reasonable for
BTC/USDT (daily volume $>\$30$B).
However, scaling to larger positions would require a market-impact model.

\textbf{Fee model.}
A flat 0.05\% maker fee is applied.
In practice, taker fees (0.10\%), funding rates, and spread costs would further
erode returns, particularly for the fast-cycle strategy that generates many fills.

\textbf{Training horizon.}
The train set (2020--2022) contains only one major bull-bear cycle.
Extended coverage of additional market regimes would strengthen the generalization
claim.

\textbf{Policy saturation.}
The learned policy outputs a near-constant $(0, 0)$ action because Bayesian
optimization drove $A_s$ to its lower bound, making the minimum action permanently
optimal.
This is a structural design limitation: co-optimizing formula coefficients and the
RL policy for the same objective removes the policy's degrees of freedom.
Future work must decouple these two optimization stages.

\textbf{$A_s$ boundary.}
The optimal $A_s = 0.101$ sits at the lower bound of the Bayesian search space,
suggesting the Bayesian stage solved the sell-gap problem completely,
leaving no room for RL to contribute further refinement.

% ─── 7. Conclusion ───────────────────────────────────────────────────────────
\section{Conclusion}
\label{sec:conclusion}

We presented a non-directional reinforcement learning framework for dynamic grid
trading on BTC/USDT in which a PPO agent controls ATR-proportional limit-order
placement without directional price prediction.
The key contributions are: (1) a five-dimensional MDP formulation encoding both
market conditions and portfolio state; (2) an eight-coefficient ATR-scaling formula
that keeps order semantics interpretable; and (3) a Bayesian search procedure that
improved validation Sharpe by $+145\%$ over hand-tuned defaults.

The final model achieves a test-set Sharpe of 43.04 — 29.2$\times$ the best
fixed-grid baseline — with a 3.12\% maximum drawdown.
Critically, post-hoc analysis shows the policy converges to a near-constant
$(0, 0)$ action output: the Bayesian optimization drove the sell-market base
coefficient $A_s$ to its lower bound, making the minimum action permanently
optimal and eliminating the need for state-conditioned selection.
The performance advantage over fixed baselines therefore stems primarily from the
Bayesian-optimized coefficient set rather than from learned adaptive behavior.

This finding reframes the research contribution: rather than demonstrating that
RL \emph{adapts} to market regimes, this work demonstrates that (a) the
ATR-proportional framework with Bayesian-optimized coefficients constitutes a
strong fixed strategy in its own right, and (b) the co-optimization of
coefficients and RL policy creates a structural conflict that prevents regime
adaptation from emerging.

\textbf{Next-phase design.}
To enable genuine regime adaptation, three changes are required:
\begin{enumerate}
    \item \textbf{Decouple coefficient and action roles.}
          Coefficients should fix the economically meaningful \emph{range}
          (e.g., sell-market in $[0.1, 1.0]\times r_t$), using Bayesian findings
          as informed bounds rather than optimized values.
          RL action then selects \emph{within} that range based on state.
    \item \textbf{Add directional state features.}
          Include rolling log-return features (e.g., 1-day and 1-week trends)
          so the policy has the signal needed to distinguish bull, bear, and
          sideways regimes.
    \item \textbf{Monitor policy entropy during training.}
          Track action standard deviation at each evaluation checkpoint;
          collapse toward zero is an early warning of the saturation failure
          mode identified here.
\end{enumerate}

Additional extensions include SAC for improved sample efficiency in the continuous
action space, and integration with the Binance live API for paper-trading
validation.

% ─── References ──────────────────────────────────────────────────────────────
\newpage
\bibliographystyle{plainnat}
\bibliography{references}

\end{document}
